% Partie : counterfactuals
% Chapitre : introduction
% Section : paradoxes-material

\documentclass[../../../include/open-logic-section]{subfiles}

\begin{document}

\olfileid[fr]{cnt}{int}{par}

\olsection{Paradoxes du conditionnel matériel}

C.~I. Lewis fut l'un des premiers à critiquer l'emploi de
$!A \lif !B$ pour symboliser les énoncés anglais de la forme
« If \dots then \dots ». Il visait notamment le conditionnel
matériel des \emph{Principia Mathematica} de Whitehead et Russell,
qui lisaient $\lif$ « implique ». Lewis objectait à juste titre
que, si $\lif$ signifiait « implique », toute proposition
fausse~$p$ impliquerait que $p$ implique~$q$, puisque
$p \lif (p \lif q)$ est vraie lorsque~$p$ est fausse ; de même,
toute proposition vraie~$q$ impliquerait que $p$ implique~$q$,
puisque $q \lif (p \lif q)$ est vraie lorsque~$q$ est vraie.

L'\emph{implication} au sens de la conséquence logique est une
relation entre des !!{formula}s ou des énoncés, et non le
connecteur du conditionnel matériel. Pour éviter la confusion, il
vaut donc mieux ne pas lire $\lif$ « implique ».\footnote{La
  lecture de « $\lif$ »
  comme « implique » reste courante chez les mathématiciens et les
  informaticiens. Les philosophes emploient parfois « seulement si »
  (en anglais, « only if ») pour éviter la confusion relevée par
  Lewis.} Dès que cette distinction est faite, l'objection précise
de Lewis disparaît : $p$ n'« implique » pas~$q$ simplement parce
que l'énoncé anglais représenté par~$p$ se trouve être faux. Pour
savoir si $p \Entails q$, il faut considérer \emph{toutes} les
!!{valuation}s ; on a $p \Entails/ q$ même si~$p$ représente un
énoncé qui, de fait, est faux.

Il reste pourtant quelque chose d'étrange dans des conditionnels de
la forme « si \dots, alors \dots », comme l'exemple de Lewis :
\begin{quote}
Si la Lune est faite de fromage vert, alors $2+2=4$.
\end{quote}
Il en va de même des inférences suivantes :
\begin{quote}
  La Lune n'est pas faite de fromage vert. Donc, si la Lune est faite
  de fromage vert, alors $2+2=4$.

  $2+2 = 4$. Donc, si la Lune est faite de fromage vert, alors
  $2+2=4$.
\end{quote}
Pourtant, si « si \dots, alors \dots » exprimait simplement
$\lif$, cet énoncé serait vrai sans difficulté et ces inférences
seraient valides sans difficulté.

Une autre difficulté vient de la tautologie $(!A \lif !B) \lor
(!B \lif !A)$. Elle suggérerait que, si l'on choisissait au hasard
deux énoncés déclaratifs~$S$ et $T$ dans un journal, l'énoncé
« Si $S$, alors $T$, ou bien si $T$, alors $S$ » devrait être vrai.

\end{document}
